\documentclass{uoc_memoria}

% Metadatos del Trabajo Final de Máster
\title{Generación acelerada en GPU de campos de distancia con signo (SDF) a partir de mallas poligonales}
\tfmsubtitle{Evaluación de técnicas de propagación y aplicación a iluminación global interactiva}
\tfmstudent{Daniel Martí Casanova}
\tfmprogram{Máster Universitario en Diseño y Programación de Videojuegos}
\tfmarea{Programación Gráfica y Motores de Videojuegos}
\tfmtutor{Gonzalo Nicolás Barreales}
\tfmpra{Joan Arnedo Moreno}
\tfmdate{11/2026}
\tfmlanguage{Castellano}
\tfmkeywords{Campos de distancia con signo, aceleración en GPU, iluminación global}

% Resumen en castellano
\tfmabstractes{%
Este Trabajo Final de Máster diseña, implementa y evalúa dos métodos de generación de campos de distancia con signo (\textit{Signed Distance Fields}, SDF) volumétricos en GPU a partir de mallas poligonales. Estas representaciones permiten consultas de distancia para trazado de rayos por software, pero su cálculo resulta costoso al aumentar la resolución del volumen y la complejidad geométrica. Se desarrolla un método base por fuerza bruta mediante \textit{compute shaders} y una alternativa de propagación espacial basada en el algoritmo \textit{Jump Flooding} (JFA) tridimensional con captura rasterizada de superficie. La validación se apoya en un visor interactivo con trazado de caminos híbrido y filtro espacio-temporal SVGF.
}

% Abstract en inglés (cumple ASD-STE100)
\tfmabstracten{%
This Master's Thesis designs, implements, and evaluates two GPU-accelerated volumetric signed distance field (SDF) generation methods from polygonal meshes. Distance fields enable efficient distance queries for software ray tracing, but computing them becomes expensive as grid resolution and geometric complexity increase. This work develops a brute-force compute shader baseline and an alternative spatial propagation pipeline using 3D Jump Flooding (JFA) with rasterized surface initialization. An interactive hybrid path tracer with SVGF spatiotemporal denoising provides the experimental validation framework.
}

\begin{document}

% -------------------------------------------------------------
% PORTADA
% -------------------------------------------------------------
\makecoverpage

% -------------------------------------------------------------
% PÁGINA DE LICENCIA
% -------------------------------------------------------------
\clearpage
\pagestyle{licensepage}

\vspace*{160mm}

\noindent
\includegraphics[height=1.0cm]{\templateassets/cc_by_nc_nd.png}\par
\vspace{2mm}
\noindent
Esta obra está sujeta a una licencia de Reconocimiento-NoComercial-SinObraDerivada 3.0 España de Creative Commons.

% -------------------------------------------------------------
% FICHA DEL TRABAJO FINAL (Página 1)
% -------------------------------------------------------------
\clearpage
\pagenumbering{arabic}
\setcounter{page}{1}
\pagestyle{uocdefault}

\makefichaproyecto

% -------------------------------------------------------------
% AGRADECIMIENTOS
% -------------------------------------------------------------
\clearpage
\section*{Agradecimientos}

Agradecimientos del autor a las personas, empresas o instituciones que han contribuido en la realización de este proyecto.

% -------------------------------------------------------------
% DECLARACIÓN DE RECURSOS DE TERCEROS
% -------------------------------------------------------------
\clearpage
\section*{Declaración de uso de recursos de terceros}

En el desarrollo de este trabajo se han utilizado bibliotecas de software de terceros, recursos multimedia externos y herramientas de Inteligencia Artificial para tareas de asistencia en edición, estructuración y verificación técnica. El registro pormenorizado se mantiene en la documentación académica del proyecto.

% -------------------------------------------------------------
% ÍNDICE Y LISTAS
% -------------------------------------------------------------
\clearpage
\section*{Índice}
\vspace{2mm}

{
  \hypersetup{linkcolor=black}
  \renewcommand{\contentsname}{}
  \vspace{-15mm}
  \tableofcontents
}

\clearpage
\section*{Lista de figuras}
{
  \hypersetup{linkcolor=black}
  \renewcommand{\listfigurename}{}
  \vspace{-15mm}
  \listoffigures
}

% -------------------------------------------------------------
% CUERPO DE LA MEMORIA (Página 1 en adelante)
% -------------------------------------------------------------
\clearpage
\pagenumbering{arabic}
\setcounter{page}{1}

\section{Introducción}

\subsection{Contexto y justificación del Trabajo}
Texto de contexto y justificación...

\subsection{Objetivos del Trabajo}
Texto de objetivos...

\subsection{Impacto en sostenibilidad, ético-social y de diversidad}
Texto de impactos...

\subsection{Enfoque y método seguido}
Texto de enfoque y metodología...

\subsection{Planificación del Trabajo}
Texto de planificación...

\subsection{Breve sumario de productos obtenidos}
Resumen de productos...

\subsection{Breve descripción de los otros capítulos de la memoria}
Descripción capitular...

\clearpage
\section{Estado del arte}
Texto del estado del arte...

\clearpage
\section{Materiales y métodos}
Texto de materiales y métodos...

\clearpage
\section{Diseño}
Texto de diseño...

\clearpage
\section{Resultados}
Texto de resultados...

\clearpage
\section{Conclusiones y Trabajo futuro}

\subsection{Conclusiones}
Texto de conclusiones...

\subsection{Seguimiento de la planificación}
Texto de seguimiento...

\subsection{Líneas de futuro}
Texto de líneas de futuro...

\clearpage
\section{Glosario}
Definición de términos...

\clearpage
\section{Bibliografía}
Lista numerada de referencias bibliográficas...

\clearpage
\section{Anexos}
Anexos complementarios...

\end{document}
